% ============================================================================
% ARBEITSBLATT DS 9 - Gruppe A (A2): Repaso y Preparación
% ============================================================================
% Spanisch Spätbeginner - Einstimmung vor Beobachtungsstunde
% Fokus: Wiederholung ser/estar, tener, hay que, Personenbeschreibung
% ============================================================================

\documentclass[11pt,a4paper]{article}

\usepackage[utf8]{inputenc}
\usepackage[T1]{fontenc}
\usepackage[ngerman]{babel}
\usepackage[left=2cm,right=2cm,top=1.8cm,bottom=1.5cm]{geometry}
\usepackage{helvet}
\usepackage[table]{xcolor}
\usepackage{tabularx}
\usepackage{array}
\usepackage{enumitem}
\usepackage{tcolorbox}
\usepackage{fancyhdr}
\usepackage{lastpage}
\usepackage{amssymb}

\renewcommand{\familydefault}{\sfdefault}

\definecolor{spanisch}{HTML}{c41e3a}
\definecolor{grau}{HTML}{666666}
\definecolor{hellgrau}{HTML}{f5f5f5}

\pagestyle{fancy}
\fancyhf{}
\fancyhead[L]{\small\textcolor{grau}{Spanisch Spätbeginner -- Gruppe A (A2)}}
\fancyhead[R]{\small\textcolor{grau}{AB-09-A}}
\fancyfoot[C]{\small\thepage/\pageref{LastPage}}
\renewcommand{\headrulewidth}{0.5pt}

\newcommand{\luecke}[1]{\underline{\hspace{#1}}}
\newcommand{\punktebox}[1]{\hfill\fbox{\small \_\_/#1}}
\newcommand{\es}[1]{\textit{#1}}

\newtcolorbox{merkkasten}[1][]{
    colback=white,
    colframe=spanisch,
    fonttitle=\bfseries\color{spanisch},
    title=#1,
    boxrule=1.5pt,
    arc=0pt,
    left=5pt, right=5pt, top=5pt, bottom=5pt
}

\newtcolorbox{hilfekasten}{
    colback=hellgrau,
    colframe=grau,
    boxrule=0.5pt,
    arc=0pt,
    left=5pt, right=5pt, top=3pt, bottom=3pt
}

\newcounter{aufgabe}
\newcommand{\aufgabe}[2]{%
    \stepcounter{aufgabe}%
    \vspace{0.8em}%
    \noindent\textbf{\theaufgabe.} #1 \punktebox{#2}%
    \vspace{0.3em}%
}

\begin{document}

% ============================================================================
% KOPFBEREICH
% ============================================================================
\noindent
\begin{tabularx}{\textwidth}{@{}Xr@{}}
    {\Large\textbf{\textcolor{spanisch}{Repaso: Describir personas}}} & Name: \luecke{5cm} \\[0.3em]
    {\small DS 9 -- Wiederholung zur Vorbereitung} & Datum: \luecke{3cm} \\
\end{tabularx}

\vspace{0.3em}
\hrule
\vspace{0.8em}

% ============================================================================
% MERKKASTEN 1: ser vs. estar
% ============================================================================
\begin{merkkasten}[Kasten 1: \es{ser} vs. \es{estar} -- Wann benutze ich was?]
\vspace{0.2em}

\begin{tabularx}{\textwidth}{@{}|p{2.5cm}|X|X|@{}}
\hline
\rowcolor{hellgrau}
\textbf{Verb} & \textbf{Verwendung} & \textbf{Beispiel} \\
\hline
\textbf{ser} & \luecke{4cm} (dauerhaft) & \es{Soy} \luecke{2.5cm} (Ich bin sympathisch.) \\[0.6em]
\hline
\textbf{estar} & \luecke{4cm} (vorübergehend) & \es{Estoy} \luecke{2.5cm} (Ich bin müde.) \\[0.6em]
\hline
\end{tabularx}

\vspace{0.5em}
\textbf{Merke:} \es{ser} = \luecke{6cm} \hspace{1cm} \es{estar} = \luecke{6cm}
\vspace{0.3em}
\end{merkkasten}

\vspace{0.8em}

% ============================================================================
% MERKKASTEN 2: hay que + Infinitiv
% ============================================================================
\begin{merkkasten}[Kasten 2: \es{hay que} + Infinitiv]
\vspace{0.2em}

\textbf{Bedeutung:} \es{hay que} = \luecke{5cm}

\vspace{0.3em}
\textbf{Struktur:} \es{hay que} + \luecke{3cm}

\vspace{0.3em}
\textbf{Beispiel:} \es{Para ser médico hay que ser} \luecke{3cm}.
\vspace{0.3em}
\end{merkkasten}

\vspace{1em}

% ============================================================================
% AUFGABEN
% ============================================================================

\aufgabe{Ergänze \es{ser} oder \es{estar} in der richtigen Form.}{6}

\begin{hilfekasten}
\textbf{Hilfe:} \es{ser} = Eigenschaften (wie jemand IST) | \es{estar} = Zustände (wie jemand sich FÜHLT)
\end{hilfekasten}

\vspace{0.3em}
\begin{enumerate}[label=\alph*), leftmargin=2em]
    \item Mi hermano \luecke{2cm} muy inteligente.
    \item Hoy \luecke{2cm} cansada porque no he dormido bien.
    \item ¿Cómo \luecke{2cm} tu mejor amiga? -- Es simpática.
    \item Mis padres \luecke{2cm} muy trabajadores.
    \item ¿Por qué \luecke{2cm} triste hoy?
    \item Nosotros \luecke{2cm} de Alemania.
\end{enumerate}

\vspace{0.5em}

\aufgabe{Wortschatz: Adjektive für Personen. Übersetze.}{8}

\begin{tabularx}{\textwidth}{@{}|X|X|X|X|@{}}
\hline
\rowcolor{hellgrau}
\textbf{Spanisch} & \textbf{Deutsch} & \textbf{Spanisch} & \textbf{Deutsch} \\
\hline
simpático/-a & \luecke{2.5cm} & paciente & \luecke{2.5cm} \\[0.5em]
\hline
inteligente & \luecke{2.5cm} & divertido/-a & \luecke{2.5cm} \\[0.5em]
\hline
trabajador/-a & \luecke{2.5cm} & responsable & \luecke{2.5cm} \\[0.5em]
\hline
creativo/-a & \luecke{2.5cm} & organizado/-a & \luecke{2.5cm} \\[0.5em]
\hline
\end{tabularx}

\vspace{1em}

\aufgabe{Was braucht man für diese Berufe? Ergänze mit \es{hay que ser...}}{6}

\begin{hilfekasten}
\textbf{Adjektive:} \es{paciente} | \es{responsable} | \es{inteligente} | \es{creativo/-a} | \es{divertido/-a} | \es{organizado/-a}
\end{hilfekasten}

\vspace{0.3em}
\begin{enumerate}[label=\alph*), leftmargin=2em]
    \item Para ser \textbf{médico/-a}, hay que ser \luecke{4cm}.
    \item Para ser \textbf{profesor/-a}, hay que ser \luecke{4cm}.
    \item Para ser \textbf{artista}, hay que ser \luecke{4cm}.
\end{enumerate}

\vspace{0.5em}

\aufgabe{Partnerinterview: Beschreibe deinen Partner / deine Partnerin.}{7}

\begin{hilfekasten}
\textbf{Fragen:} \es{¿Cómo te llamas? ¿Cuántos años tienes? ¿Cómo eres? ¿Qué te gusta?}
\end{hilfekasten}

\vspace{0.3em}
\textbf{Notiere 4 Informationen über deinen Partner / deine Partnerin:}

\vspace{0.3cm}
1. \luecke{13cm}

\vspace{0.4cm}
2. \luecke{13cm}

\vspace{0.4cm}
3. \luecke{13cm}

\vspace{0.4cm}
4. \luecke{13cm}

\vspace{1em}

\aufgabe{Schreibe einen kurzen Text: Beschreibe eine Person (4--5 Sätze).}{8}

\begin{hilfekasten}
\textbf{Verwende:} \es{ser + Adjektiv}, \es{tener + años}, \es{también}, \es{porque}
\end{hilfekasten}

\vspace{0.5em}
\textbf{Mi compañero/-a ideal / Mi mejor amigo/-a}

\vspace{0.3em}
\begin{tabularx}{\textwidth}{@{}|X|@{}}
\hline
\\[0.5em]
\luecke{14cm} \\[0.8em]
\luecke{14cm} \\[0.8em]
\luecke{14cm} \\[0.8em]
\luecke{14cm} \\[0.8em]
\luecke{14cm} \\[0.8em]
\hline
\end{tabularx}

\vspace{1em}
\hrule
\vspace{0.3em}
\begin{flushright}
\textbf{Gesamt: \luecke{1cm} / 35 Punkte}
\end{flushright}

\end{document}
